% !TeX program = xelatex
% File này có thể biên dịch độc lập trên Overleaf bằng XeLaTeX.
% Nếu báo cáo đã có file main.tex, chỉ sao chép phần từ
% \chapter{ĐẶT VẤN ĐỀ} đến trước \end{document}.

\documentclass[12pt,a4paper,oneside]{report}

\usepackage{iftex}
\ifPDFTeX
  \usepackage[utf8]{inputenc}
  \usepackage[T5]{fontenc}
  \usepackage[vietnamese]{babel}
\else
  \usepackage{fontspec}
  \setmainfont{TeX Gyre Termes}
  \usepackage[vietnamese]{babel}
\fi

\usepackage[
  a4paper,
  top=2.5cm,
  bottom=2.5cm,
  left=3.0cm,
  right=2.0cm
]{geometry}
\usepackage{setspace}
\usepackage{enumitem}
\usepackage{booktabs}
\usepackage{tabularx}
\usepackage{array}
\usepackage{xcolor}
\usepackage[hidelinks]{hyperref}

\onehalfspacing
\setlength{\parindent}{1.25cm}
\setlength{\parskip}{0.25em}
\setcounter{secnumdepth}{3}
\renewcommand{\arraystretch}{1.25}
\newcolumntype{Y}{>{\raggedright\arraybackslash}X}

\begin{document}

\chapter{ĐẶT VẤN ĐỀ}
\label{chap:dat-van-de}

\section{Bối cảnh của bài toán}
\label{sec:boi-canh}

Trong quá trình sản xuất, không phải toàn bộ nguyên vật liệu đầu vào đều được
chuyển hóa thành sản phẩm cuối cùng. Nhiều doanh nghiệp thường xuyên phát sinh
vật liệu dư thừa, phụ phẩm hoặc phế liệu vẫn còn giá trị sử dụng. Chẳng hạn, một
nhà máy nhựa có thể phát sinh khoảng 20 tấn nhựa PP mỗi tháng; một xưởng gỗ có
hàng nghìn pallet không còn nhu cầu sử dụng; một nhà máy thực phẩm có bã cà phê;
trong khi một doanh nghiệp dệt may có lượng lớn vải vụn sau quá trình cắt may.

Ở chiều ngược lại, nhiều doanh nghiệp khác vẫn phải mua nguyên liệu mới có tính
chất tương tự để phục vụ hoạt động sản xuất. Một nhà máy có thể đang cần nhựa PP
tái sử dụng, một cơ sở đóng gói có thể cần pallet gỗ, hoặc một đơn vị sản xuất
vật liệu sinh học có thể sử dụng bã cà phê làm đầu vào. Như vậy, nguồn vật liệu
dư của doanh nghiệp này có khả năng trở thành nguyên liệu đầu vào của doanh
nghiệp khác.

Mối quan hệ trên là một biểu hiện của kinh tế tuần hoàn, trong đó vòng đời của
vật liệu được kéo dài thông qua hoạt động tái sử dụng, tái chế hoặc chuyển giao
cho một quy trình sản xuất khác. Cách tiếp cận này vừa giúp bên cung giảm chi phí
xử lý chất thải, vừa giúp bên mua tiếp cận nguồn nguyên liệu thay thế với chi phí
hợp lý hơn, đồng thời góp phần giảm lượng chất thải đưa ra môi trường.

Tuy nhiên, việc hình thành các giao dịch vật liệu tuần hoàn giữa doanh nghiệp
hiện vẫn gặp nhiều trở ngại. Thông tin về nguồn cung và nhu cầu thường nằm rời
rạc trong từng doanh nghiệp, được trao đổi qua quan hệ cá nhân, điện thoại, thư
điện tử hoặc các nhóm mạng xã hội. Do thiếu một nền tảng B2B chuyên biệt, một
nguồn vật liệu còn giá trị có thể bị xử lý như chất thải chỉ vì doanh nghiệp sở
hữu không tìm được bên có nhu cầu tại đúng thời điểm.

Từ bối cảnh đó, đề tài đề xuất xây dựng \textbf{Circular Materials Exchange}
-- nền tảng trao đổi vật liệu tuần hoàn giữa các doanh nghiệp. Hệ thống đóng vai
trò là không gian số tập trung để công bố nguồn cung, công bố nhu cầu, tìm kiếm
đối tác, trao đổi đề nghị và theo dõi quá trình giao dịch.

\section{Hiện trạng và các vấn đề cần giải quyết}
\label{sec:hien-trang}

\subsection{Thông tin cung và cầu bị phân tán}

Doanh nghiệp có vật liệu dư thừa chưa có một kênh thống nhất để mô tả loại vật
liệu, số lượng, chất lượng, vị trí, đơn giá và điều kiện giao nhận. Tương tự,
doanh nghiệp cần mua vật liệu cũng thiếu nơi công bố nhu cầu theo một cấu trúc
chuẩn. Việc thiếu dữ liệu tập trung khiến hai phía khó nhận biết nhau, ngay cả
khi nguồn cung và nhu cầu có khả năng phù hợp.

\subsection{Khó tìm kiếm và đánh giá cơ hội giao dịch}

Khi thông tin được trao đổi qua nhiều kênh khác nhau, doanh nghiệp phải mất thời
gian tổng hợp và xác minh dữ liệu. Người mua khó lọc vật liệu theo nhóm, vị trí,
số lượng hoặc mức giá. Người bán cũng khó nhận biết doanh nghiệp nào đang có nhu
cầu phù hợp. Thông tin về tư cách doanh nghiệp và lịch sử đánh giá đối tác chưa
được quản lý một cách có hệ thống, làm tăng rủi ro khi thiết lập quan hệ giao
dịch mới.

\subsection{Quy trình thương lượng và theo dõi chưa được số hóa}

Sau khi tìm thấy một nguồn vật liệu phù hợp, các bước gửi đề nghị, chấp nhận hoặc
từ chối đề nghị, ghi nhận thỏa thuận và cập nhật tiến độ thường tiếp tục được
thực hiện thủ công. Không có một mã giao dịch và lịch sử trạng thái thống nhất
để cả bên mua, bên bán và người quản trị cùng theo dõi. Điều này có thể gây khó
khăn khi cần xác định trách nhiệm, đối chiếu nội dung thỏa thuận hoặc đánh giá kết
quả giao dịch.

\subsection{Thiếu cơ chế tạo dựng niềm tin}

Giao dịch B2B liên quan đến thông tin pháp lý của doanh nghiệp, chất lượng vật
liệu và khả năng thực hiện cam kết. Nếu không có quy trình đăng ký hồ sơ doanh
nghiệp, kiểm duyệt và đánh giá sau giao dịch, người dùng khó xác định mức độ uy
tín của đối tác. Bên cạnh đó, việc không có thông báo tập trung khiến người dùng
có thể bỏ lỡ đề nghị hoặc thay đổi trạng thái quan trọng.

\subsection{Thanh toán và logistics chưa có hạ tầng tích hợp}

Hệ sinh thái hiện tại chưa có cơ chế logistics chuyên biệt cho nhiều loại vật
liệu công nghiệp và chưa có quy trình thanh toán điện tử được tích hợp thống
nhất. Đây là bài toán lớn, liên quan đến ngân hàng, vận chuyển, kiểm định và pháp
lý. Trong phạm vi prototype của môn học, hệ thống chỉ mô phỏng việc ghi nhận
thanh toán, phí giao dịch, escrow và trạng thái giao nhận; việc chuyển tiền và
vận chuyển thực tế vẫn được các doanh nghiệp thực hiện ngoài hệ thống.

\subsection{Hệ quả của hiện trạng}

Các hạn chế nêu trên dẫn đến ba nhóm hệ quả chính:

\begin{itemize}[leftmargin=1.1cm]
  \item \textbf{Về kinh tế:} doanh nghiệp phát sinh thêm chi phí xử lý vật liệu
  dư, trong khi doanh nghiệp khác vẫn phải mua nguyên liệu mới với chi phí cao.

  \item \textbf{Về vận hành:} quá trình tìm đối tác, thương lượng và theo dõi
  giao dịch phụ thuộc nhiều vào thao tác thủ công, khó kiểm soát và khó truy vết.

  \item \textbf{Về môi trường:} vật liệu còn khả năng tái sử dụng hoặc tái chế
  vẫn bị đưa vào luồng chất thải, làm giảm hiệu quả sử dụng tài nguyên và tăng áp
  lực lên hoạt động xử lý chất thải.
\end{itemize}

\section{Nhu cầu xây dựng hệ thống}
\label{sec:nhu-cau}

Từ các vấn đề trên, cần xây dựng một hệ thống thông tin có khả năng chuẩn hóa và
kết nối các hoạt động chính trong quá trình trao đổi vật liệu tuần hoàn. Hệ thống
cần đáp ứng các nhu cầu sau:

\begin{enumerate}[leftmargin=1.1cm]
  \item Tạo một nơi công bố tập trung cho nguồn vật liệu dư thừa và nhu cầu mua
  vật liệu của doanh nghiệp.

  \item Cho phép người dùng tìm kiếm, lọc và xem thông tin chi tiết để phát hiện
  cơ hội kết nối cung -- cầu.

  \item Hỗ trợ doanh nghiệp gửi đề nghị với số lượng, mức giá và lời nhắn cụ thể;
  đồng thời cho phép bên nhận chấp nhận hoặc từ chối đề nghị.

  \item Chuyển một đề nghị đã được chấp nhận thành giao dịch có mã định danh,
  trạng thái và lịch sử sự kiện rõ ràng.

  \item Hỗ trợ hồ sơ doanh nghiệp, quy trình phê duyệt và đánh giá sau giao dịch
  nhằm tăng tính minh bạch và mức độ tin cậy giữa các bên.

  \item Cung cấp chức năng quản trị để kiểm soát doanh nghiệp, danh mục vật liệu,
  tin đăng, giao dịch và các dữ liệu tài chính ở mức prototype.

  \item Tạo nền tảng dữ liệu và kiến trúc kỹ thuật có thể mở rộng trong tương lai
  cho thuật toán gợi ý, tự động ghép cặp cung -- cầu, thanh toán điện tử và tích
  hợp logistics.
\end{enumerate}

\section{Mục tiêu và phạm vi của đề tài}
\label{sec:muc-tieu-pham-vi}

\subsection{Mục tiêu tổng quát}

Mục tiêu của đề tài là phân tích, thiết kế và xây dựng prototype một nền tảng B2B
giúp doanh nghiệp trao đổi vật liệu dư thừa theo hướng kinh tế tuần hoàn. Hệ thống
phải thể hiện được luồng xuyên suốt từ công bố vật liệu, tìm kiếm, gửi đề nghị,
hình thành giao dịch, theo dõi trạng thái đến đánh giá đối tác.

\subsection{Mục tiêu cụ thể}

\begin{itemize}[leftmargin=1.1cm]
  \item Số hóa thông tin nguồn cung và nhu cầu vật liệu theo cấu trúc thống nhất.
  \item Rút ngắn quá trình tìm kiếm và liên hệ giữa doanh nghiệp mua và bán.
  \item Ghi nhận đầy đủ đề nghị, giao dịch và lịch sử thay đổi trạng thái.
  \item Hỗ trợ kiểm duyệt doanh nghiệp và quản trị dữ liệu hệ thống.
  \item Xây dựng backend theo kiến trúc microservice để phân tách trách nhiệm
  nghiệp vụ và tạo khả năng mở rộng.
  \item Triển khai được một phiên bản hoạt động trên môi trường máy chủ để phục
  vụ thử nghiệm và trình diễn.
\end{itemize}

\subsection{Phạm vi MVP}

Trong phạm vi bài tập lớn, sản phẩm được xác định là một MVP/prototype học thuật.
Các chức năng thuộc phạm vi thực hiện gồm:

\begin{itemize}[leftmargin=1.1cm]
  \item đăng ký, đăng nhập và quản lý phiên bằng JWT bearer token; middleware
  vẫn giữ tương thích với demo token cũ trong môi trường trình diễn;
  \item tạo hồ sơ doanh nghiệp và thực hiện phê duyệt hoặc từ chối bởi Admin;
  \item quản lý danh mục và nguồn cung; hiển thị nhu cầu vật liệu, trong khi thao
  tác tạo nhu cầu hiện mới ở mức stub và chưa ghi cơ sở dữ liệu;
  \item tìm kiếm, lọc và xem chi tiết tin đăng;
  \item gửi, tiếp nhận, chấp nhận hoặc từ chối đề nghị mua;
  \item tạo giao dịch, cập nhật trạng thái và lưu lịch sử giao dịch;
  \item tải ảnh vật liệu, gửi đánh giá và nhận thông báo trong ứng dụng;
  \item ghi nhận phí, escrow và ví người bán ở mức prototype;
  \item cung cấp các màn hình quản trị và thống kê cơ bản.
\end{itemize}

Các nội dung ngoài phạm vi MVP gồm tích hợp ngân hàng hoặc cổng thanh toán thật,
điều phối logistics, theo dõi vận chuyển theo thời gian thực, hợp đồng điện tử,
chữ ký số, kiểm định chất lượng bởi bên thứ ba, xác minh pháp lý tự động và quy
trình giải quyết tranh chấp hoàn chỉnh.

Mục tiêu ``tự động kết nối cung -- cầu'' là định hướng phát triển của hệ thống.
Phiên bản hiện tại chưa cài đặt thuật toán tự động matching hoặc recommendation.
Việc kết nối đang được thực hiện thông qua công bố tập trung, tìm kiếm, lọc, xem
chi tiết và gửi Offer. Cách xác định phạm vi này giúp báo cáo phản ánh đúng năng
lực của mã nguồn hiện có, đồng thời vẫn giữ được định hướng mở rộng của đề tài.

\section{Mô tả sơ bộ các yêu cầu chức năng}
\label{sec:yeu-cau-chuc-nang-so-bo}

Các tác nhân chính của hệ thống gồm Guest, Business User và Admin. Buyer và
Seller là hai vai trò nghiệp vụ của Business User, không phải hai loại tài khoản
tách biệt; một doanh nghiệp có thể vừa mua vừa bán tùy từng giao dịch.

\begin{table}[htbp]
  \centering
  \caption{Các tác nhân chính của hệ thống}
  \label{tab:tac-nhan-chinh}
  \begin{tabularx}{\textwidth}{p{3.0cm}Y}
    \toprule
    \textbf{Tác nhân} & \textbf{Vai trò sơ bộ} \\
    \midrule
    Guest &
    Xem marketplace, danh sách nhu cầu và chi tiết vật liệu; tìm kiếm thông tin;
    đăng ký và đăng nhập. \\

    Business User &
    Quản lý hồ sơ cá nhân, tạo hồ sơ doanh nghiệp, theo dõi trạng thái phê duyệt
    và sử dụng các chức năng nghiệp vụ sau khi đăng nhập. \\

    Buyer &
    Đăng nhu cầu, tìm nguồn cung, gửi đề nghị mua, theo dõi giao dịch và đánh giá
    Seller. \\

    Seller &
    Đăng nguồn cung, tiếp nhận và xử lý đề nghị, cập nhật giao dịch, theo dõi ví
    và đánh giá Buyer. \\

    Admin &
    Duyệt doanh nghiệp; quản lý danh mục, tin đăng và giao dịch; theo dõi dữ liệu
    tài chính, escrow và các yêu cầu quản trị. \\
    \bottomrule
  \end{tabularx}
\end{table}

Từ góc nhìn chức năng, hệ thống được chia sơ bộ thành các nhóm sau:

\begin{table}[htbp]
  \centering
  \caption{Các nhóm yêu cầu chức năng sơ bộ}
  \label{tab:nhom-chuc-nang}
  \begin{tabularx}{\textwidth}{p{3.5cm}Y}
    \toprule
    \textbf{Nhóm chức năng} & \textbf{Mô tả} \\
    \midrule
    Tài khoản và xác thực &
    Đăng ký, đăng nhập, phát hành JWT bearer token, lấy thông tin người dùng hiện
    tại và bảo vệ các chức năng yêu cầu xác thực. Middleware vẫn hỗ trợ token
    demo cũ để không làm gián đoạn phiên trình diễn. \\

    Hồ sơ doanh nghiệp &
    Tạo và xem hồ sơ doanh nghiệp; Admin duyệt hoặc từ chối; chỉ doanh nghiệp đã
    được xác minh mới được thực hiện các thao tác mua hoặc đăng bán quan trọng. \\

    Quản lý vật liệu &
    Xem danh mục; đăng và xem nguồn cung; lưu thông tin mô tả, số lượng, đơn vị,
    mức giá, vị trí, đóng gói và hình ảnh. Route cập nhật/xóa nguồn cung và tạo
    nhu cầu hiện mới ở mức stub. \\

    Khám phá cung -- cầu &
    Công khai marketplace; tìm kiếm và lọc nguồn cung hoặc nhu cầu; xem chi tiết
    để người dùng tự xác định cơ hội phù hợp trong phạm vi MVP. \\

    Đề nghị giao dịch &
    Buyer gửi Offer gắn với nguồn cung, trong đó có số lượng, giá đề xuất và lời
    nhắn; Seller xem, chấp nhận hoặc từ chối đề nghị. \\

    Giao dịch và timeline &
    Tự động tạo Transaction khi Offer được chấp nhận; cho phép xem danh sách,
    chi tiết, cập nhật trạng thái và lưu TransactionEvent để hình thành lịch sử
    xử lý. \\

    Tài chính prototype &
    Ghi nhận số tiền, phí nền tảng, escrow, giải ngân, ví Seller và yêu cầu rút
    tiền. Đây là sổ ghi nhận trong prototype, không phải tích hợp chuyển tiền
    ngân hàng thực tế. \\

    Đánh giá và thông báo &
    Tạo đánh giá sau giao dịch; xem điểm và nhận xét; nhận thông báo, đánh dấu đã
    đọc hoặc đánh dấu tất cả đã đọc. \\

    Quản trị hệ thống &
    Cung cấp dashboard và các chức năng quản lý doanh nghiệp, danh mục, tin đăng,
    giao dịch, báo cáo, tài chính và escrow. \\
    \bottomrule
  \end{tabularx}
\end{table}

Luồng nghiệp vụ cốt lõi của hệ thống được mô tả sơ bộ như sau:

\begin{enumerate}[leftmargin=1.1cm]
  \item Người dùng đăng ký, đăng nhập và tạo hồ sơ doanh nghiệp.
  \item Admin xem xét, sau đó phê duyệt hoặc từ chối hồ sơ.
  \item Seller đăng nguồn cung vật liệu; Buyer xem nhu cầu đã công bố. Chức năng
  tạo nhu cầu hiện mới là stub và chưa ghi cơ sở dữ liệu.
  \item Buyer tìm kiếm marketplace, xem chi tiết nguồn cung và gửi Offer.
  \item Seller tiếp nhận Offer và quyết định chấp nhận hoặc từ chối.
  \item Khi Offer được chấp nhận, hệ thống tạo Transaction và sự kiện khởi tạo.
  \item Hai bên tự thỏa thuận thanh toán và giao nhận bên ngoài nền tảng; hệ thống
  chỉ ghi nhận trạng thái và timeline tương ứng.
  \item Sau khi hoàn tất, các bên có thể đánh giá nhau; dữ liệu được dùng để tăng
  tính minh bạch cho những giao dịch tiếp theo.
\end{enumerate}

\section{Mô tả sơ bộ giải pháp xây dựng hệ thống}
\label{sec:giai-phap-so-bo}

\subsection{Giải pháp nghiệp vụ}

Giải pháp được xây dựng theo mô hình sàn giao dịch B2B. Thay vì trực tiếp sở hữu
vật liệu hoặc thực hiện vận chuyển, hệ thống tập trung vào quản lý thông tin và
quy trình kết nối. Mỗi nguồn cung và nhu cầu được biểu diễn bằng một bản ghi có
cấu trúc. Người dùng có thể tìm kiếm dữ liệu, gửi đề nghị và theo dõi giao dịch
trên cùng một nền tảng.

Hồ sơ doanh nghiệp và trạng thái phê duyệt được sử dụng như một lớp kiểm soát
trước khi người dùng thực hiện các thao tác quan trọng. Offer đóng vai trò là đối
tượng thương lượng giữa Buyer và Seller. Khi Offer được chấp nhận, hệ thống tạo
Transaction và TransactionEvent để tách giai đoạn thương lượng khỏi giai đoạn
thực hiện giao dịch. Review và Notification hỗ trợ tạo dựng uy tín và duy trì
thông tin liên tục cho người dùng.

\subsection{Giải pháp kiến trúc}

Backend được tổ chức theo định hướng microservice. Frontend gửi yêu cầu REST đến
API Gateway; Gateway thực hiện middleware xác thực, kiểm tra quyền và chuyển đổi
REST/JSON thành lời gọi gRPC. Sáu service nghiệp vụ xử lý theo cấu trúc
Handler--Service--Repository và chỉ repository của service mới truy cập cơ sở dữ
liệu thuộc miền đó. Gateway không chứa SQL hoặc driver PostgreSQL.

Hệ thống hiện gồm sáu microservice chính:

\begin{table}[htbp]
  \centering
  \caption{Các microservice của hệ thống}
  \label{tab:danh-sach-microservice}
  \begin{tabularx}{\textwidth}{p{3.4cm}p{2.6cm}Y}
    \toprule
    \textbf{Thành phần} & \textbf{Cơ sở dữ liệu} & \textbf{Trách nhiệm} \\
    \midrule
    Auth Service & \texttt{auth\_db} &
    Đăng ký, đăng nhập, xác thực và thông tin người dùng. \\

    Company Service & \texttt{company\_db} &
    Hồ sơ doanh nghiệp và quy trình phê duyệt. \\

    Material Service & \texttt{material\_db} &
    Danh mục, nguồn cung, nhu cầu, tìm kiếm và lọc vật liệu. \\

    Order Service & \texttt{order\_db} &
    Offer, Transaction, TransactionEvent và dữ liệu tài chính liên quan. \\

    Review Service & \texttt{review\_db} &
    Đánh giá, nhận xét và điểm trung bình của đối tác. \\

    Notification Service & \texttt{notif\_db} &
    Thông báo trong ứng dụng và trạng thái đã đọc. \\
    \bottomrule
  \end{tabularx}
\end{table}

Theo thiết kế, mỗi service sở hữu một cơ sở dữ liệu logic riêng theo nguyên tắc
\textit{Database per Service}. Trong môi trường project, sáu cơ sở dữ liệu được
đặt trong cùng một container PostgreSQL 15 để giảm chi phí triển khai. Các quan
hệ xuyên database được lưu bằng ID logic. Khi một luồng cần dữ liệu từ nhiều
miền, Gateway điều phối các lời gọi gRPC, ví dụ kiểm tra Company trước khi tạo
Listing hoặc tra User trước khi tạo Review.

\subsection{Giải pháp công nghệ}

Các công nghệ chính được lựa chọn gồm:

\begin{itemize}[leftmargin=1.1cm]
  \item \textbf{Frontend:} React, TypeScript và Vite để xây dựng ứng dụng web
  dạng SPA; Nginx phục vụ tập tin tĩnh và chuyển tiếp các yêu cầu
  \texttt{/api/}.

  \item \textbf{API Gateway và microservice:} ngôn ngữ Go; Gin cho REST API tại
  Gateway; gRPC và Protocol Buffers cho giao tiếp đồng bộ giữa các service.

  \item \textbf{Xác thực và phân quyền:} Auth Service phát JWT; middleware gọi
  Auth Service để xác minh token và kiểm tra quyền Admin.

  \item \textbf{Lưu trữ dữ liệu:} PostgreSQL 15 với sáu database logic; MinIO
  được cấu hình làm kho đối tượng cho tệp và hình ảnh.

  \item \textbf{Giao tiếp bất đồng bộ:} Order Service công bố sự kiện
  \texttt{cme.orders.*} qua NATS; Notification Service subscribe theo queue group
  và dùng mã tham chiếu để chống ghi thông báo trùng.

  \item \textbf{Đóng gói và triển khai:} Docker và Docker Compose cho backend;
  frontend và backend được triển khai trên hai máy chủ tách biệt trong môi
  trường demo.
\end{itemize}

Kiến trúc trên phù hợp với mục tiêu của bài tập vì giúp thể hiện rõ ranh giới
nghiệp vụ, lớp giao diện, lớp dịch vụ và lớp lưu trữ. Đồng thời, việc tách service
tạo điều kiện mở rộng từng nhóm chức năng mà không phải chuyển toàn bộ hệ thống
thành một khối nguyên khối.

\subsection{Định hướng mở rộng}

Từ nền tảng dữ liệu của MVP, hệ thống có thể phát triển thêm các khả năng sau:

\begin{itemize}[leftmargin=1.1cm]
  \item xây dựng thuật toán matching dựa trên loại vật liệu, thuộc tính kỹ thuật,
  số lượng, vị trí, thời hạn và mức giá;
  \item gợi ý nguồn cung hoặc nhu cầu phù hợp theo lịch sử tương tác;
  \item tích hợp cổng thanh toán và dịch vụ ngân hàng để thay thế cơ chế tài
  chính mô phỏng;
  \item tích hợp doanh nghiệp logistics và theo dõi vận chuyển;
  \item bổ sung kiểm định chất lượng, hợp đồng điện tử và xử lý tranh chấp;
  \item bổ sung retry/dead-letter cho sự kiện và cơ chế giám sát, tracing
  microservice.
\end{itemize}

\section{Kết chương}
\label{sec:ket-chuong-1}

Chương này đã trình bày bối cảnh phát sinh vật liệu dư thừa trong sản xuất, những
hạn chế của phương thức kết nối hiện tại và nhu cầu xây dựng một nền tảng trao
đổi vật liệu tuần hoàn B2B. Circular Materials Exchange được đề xuất nhằm tập
trung hóa thông tin cung -- cầu, số hóa quá trình gửi đề nghị và theo dõi giao
dịch, đồng thời hỗ trợ kiểm duyệt và đánh giá để tăng mức độ tin cậy.

Trong phạm vi bài tập lớn, hệ thống được triển khai ở mức MVP với kiến trúc
microservice, gồm API Gateway và sáu service nghiệp vụ. Các chức năng thanh toán,
escrow và giao nhận chỉ được ghi nhận hoặc mô phỏng; thuật toán tự động ghép cặp
cung -- cầu được xác định là hướng phát triển tiếp theo. Các chương sau sẽ tiếp
tục phân tích nghiệp vụ, đặc tả yêu cầu và trình bày thiết kế chi tiết cho hệ
thống.

\end{document}
